% !PRINTED 368
\weberpage{368}
Nach dem oben Bemerkten giebt es aber auch immer Lösungen der Gleichung (4), wenn $\gamma$ durch den grössten gemeinschaftlichen Theiler von $\alpha$ und $\beta$ theilbar ist. Insbesondere können wir den Satz aussprechen: \gesperrt{Der grösste gemeinschaftliche Theiler $\delta$ zweier Zahlen $\alpha$, $\beta$ lässt sich in der Form}
\[
\alpha x - \beta y = \delta
\]
\gesperrt{darstellen, worin $x$, $y$ ganze Zahlen sind.}

Das Theorem III lässt sich in folgender Weise verallgemeinern:

\gesperrt{Sind $\alpha_1$, $\alpha_2$, $\alpha_3 \ldots$ gegebene ganze Zahlen in beliebiger Anzahl, ohne einen allen gemeinsamen Theiler, $m$ eine beliebig gegebene Zahl, so lassen sich immer die ganzen Zahlen $x_1$, $x_2$, $x_3 \ldots$ so bestimmen, dass}

\weq{5}{m = \alpha_1 x_1 + \alpha_2 x_2 + \alpha_3 x_3 + \cdots}
\gesperrt{wird.}

Besteht die rechte Seite von (5) nur aus zwei Gliedern, so fällt dieser Satz mit dem Theorem III zusammen. Nehmen wir also die Möglichkeit, eine Formel (5) zu befriedigen, als bewiesen an, wenn die rechte Seite aus weniger Gliedern besteht, so können wir, wenn $\delta$ der grösste gemeinschaftliche Theiler von $\alpha_2$, $\alpha_3 \ldots$ ist, da $\delta$ relativ prim zu $\alpha_1$ sein muss, die beiden Gleichungen:
\[
\delta = \alpha_2 y_2 + \alpha_3 y_3 + \cdots
\]
\[
m = \alpha_1 x_1 + \delta x
\]
durch ganzzahlige $x$, $x_1$, $y_2$, $y_3 \ldots$ befriedigen. Dann ist aber die Gleichung (5) durch
\[
x_1, \quad x_2 = x y_2, \quad x_3 = x y_3 \ldots
\]
befriedigt.

Das durch den Satz III gelöste Problem wird in der Zahlentheorie auch so ausgedrückt: Sind $\alpha$, $\beta$, $\gamma$ gegebene Zahlen, so soll eine Zahl $y$ gefunden werden, die der Congruenz
\[
\alpha y \equiv \gamma \ (\mathrm{mod}\ \beta)
\]
genügt. Diese Aufgabe hat also immer eine Lösung, wenn $\gamma$ durch den grössten gemeinschaftlichen Theiler von $\alpha$ und $\beta$ theilbar ist.
